\section{Discussion}
\label{sec:discussion}

\paragraph{Bridging the open-source gap.}
\baseline~\cite{cropper2025} was primarily developed and evaluated on
Gemini 1.5 Pro~\cite{gemini2024}, a large closed-source VLM with
substantially greater capacity than open-source alternatives.  Most of
the paper's design decisions---retrieval strategy, single-temperature
generation, scorer composition, and selection strategy---were tuned on
Gemini and transferred to Mantis without modification.  A more capable
VLM like Gemini can compensate for suboptimal pipeline choices through
sheer model capacity: it may produce diverse candidates even at low
temperature, or generate meaningful crops despite redundant ICL
examples.  Mantis, being a smaller 8B-parameter model, lacks this
compensating capacity.  Our contributions---diverse retrieval,
multi-temperature sampling, expanded candidate space, and scorer
debiasing---specifically address the bottlenecks that emerge when using
a less capable, open-source VLM.  We believe this makes our
improvements particularly relevant for practical deployment, where
open-source models are often preferred for cost, privacy, and
reproducibility reasons.

\paragraph{What transfers across VLMs and what does not.}
Our experiments reveal that several design choices validated on Gemini
do not hold for Mantis.  Final-iteration selection, which improves
Gemini by +0.034 IoU, has no effect on Mantis.  Adding a single
instruction sentence to the prompt causes a -0.041 IoU drop on
Mantis---a level of prompt sensitivity unlikely to occur with a larger
model.  This highlights a broader concern for VLM-based vision
pipelines: \emph{design choices must be validated per-model}, especially
when transferring from closed-source to open-source backbones.

\paragraph{The refinement paradox.}
Iterative refinement is presented as a core contribution of
\baseline~\cite{cropper2025}: by feeding scorer outputs back to the
VLM, the system should progressively improve its crops.  In practice,
we observe that the VLM frequently produces unparseable output during
refinement iterations, and that increasing iterations from $L{=}2$ to
$L{=}4$ yields only marginal gains.  This suggests that for smaller
VLMs, single-pass candidate generation with better diversity
(multi-temperature, larger $R$) may be more effective than iterative
refinement.

\paragraph{Limitations.}
Our improvements reduce the large-crop bias but do not eliminate it:
the VLM still tends toward larger crops, especially on images where the
GT crop is a focused sub-region.  Our modifications have been validated
only on Mantis-8B-Idefics2; we expect similar or larger gains on other
open-source VLMs of comparable scale, but this remains to be tested.
